\documentclass[a4paper,11pt]{article}

% Polskie znaki i kodowanie
\usepackage[utf8]{inputenc}
\usepackage[T1]{fontenc}
\usepackage[polish]{babel}

% Pakiety do formatowania
\usepackage{geometry}
\usepackage{xcolor}
\usepackage{listings}
\usepackage{hyperref}
\usepackage{graphicx}
\usepackage{float}
\usepackage{titlesec}
\usepackage{booktabs} % Do ładniejszych tabel

% Ustawienia strony
\geometry{
 a4paper,
 total={170mm,257mm},
 left=20mm,
 top=20mm,
}

% Kolory do listingów kodu
\definecolor{codegreen}{rgb}{0,0.6,0}
\definecolor{codegray}{rgb}{0.5,0.5,0.5}
\definecolor{codepurple}{rgb}{0.58,0,0.82}
\definecolor{backcolour}{rgb}{0.96,0.96,0.96}

% Styl kodu CMake
\lstdefinestyle{cmakestyle}{
    backgroundcolor=\color{backcolour},   
    commentstyle=\color{codegreen},
    keywordstyle=\color{blue},
    numberstyle=\tiny\color{codegray},
    stringstyle=\color{codepurple},
    basicstyle=\ttfamily\footnotesize,
    breakatwhitespace=false,         
    breaklines=true,                 
    captionpos=b,                    
    keepspaces=true,                 
    numbers=left,                    
    numbersep=5pt,                  
    showspaces=false,                
    showstringspaces=false,
    showtabs=false,                  
    tabsize=2,
    language=make
}

% Styl kodu C++
\lstdefinestyle{cppstyle}{
    backgroundcolor=\color{backcolour},   
    commentstyle=\color{codegreen},
    keywordstyle=\color{magenta},
    numberstyle=\tiny\color{codegray},
    stringstyle=\color{codepurple},
    basicstyle=\ttfamily\footnotesize,
    breakatwhitespace=false,         
    breaklines=true,                 
    captionpos=b,                    
    keepspaces=true,                 
    numbers=left,                    
    numbersep=5pt,                  
    showspaces=false,                
    showstringspaces=false,
    showtabs=false,                  
    tabsize=2,
    language=C++
}

% Styl kodu Bash
\lstdefinestyle{bashstyle}{
    backgroundcolor=\color{backcolour},   
    basicstyle=\ttfamily\footnotesize,
    breaklines=true,                 
    captionpos=b,
    frame=single
}

\lstset{style=cppstyle}

% Metadane dokumentu
\title{\textbf{Dokumentacja Biblioteki PUTM EV CAN}}
\author{Michał Błotniak}
\date{\today}

\begin{document}

\maketitle

\begin{abstract}
Niniejszy dokument stanowi instrukcję obsługi biblioteki \texttt{PUTM\_EV\_CAN\_LIBRARY}. Biblioteka została zaprojektowana dla systemów wbudowanych (STM32) oraz środowiska ROS2. Główne cechy to: statyczna alokacja pamięci, obsługa wielu rodzin STM32 (F4, G4, L4, H7) dzięki warstwie abstrakcji HAL, odbiór danych oparty na przerwaniach (Interrupt-Driven) oraz wbudowana diagnostyka LED.
\end{abstract}

\tableofcontents
\newpage

\section{Wymagania systemowe i zależności}

Prawidłowa kompilacja i działanie biblioteki \texttt{PUTM\_EV\_CAN\_LIBRARY} wymaga zapewnienia odpowiedniego środowiska programistycznego. Poniższa tabela przedstawia minimalne wersje narzędzi oraz bibliotek wykorzystywanych w projekcie.

\begin{table}[h!]
\centering
\begin{tabular}{|l|l|p{6cm}|}
\hline
\textbf{Narzędzie / Biblioteka} & \textbf{Wersja Minimalna} & \textbf{Uwagi} \\ \hline
CMake & 3.16 & Wymagane przez konfigurację projektu. \\ \hline
Standard C++ & C++20 & Wymuszone w pliku \texttt{CMakeLists.txt}. \\ \hline
Standard C & C11 & Standard dla kodu C. \\ \hline
Python & 3.8+ & Niezbędny do uruchomienia skryptów generujących kod z DBC. \\ \hline
Kompilator C++ (ARM) & GCC $\ge$ 10.0 & Musi obsługiwać flagi standardu C++20 (np. \texttt{arm-none-eabi-gcc}). \\ \hline
Biblioteka HAL & STM32XXxx HAL & Zależność sprzętowa dla mikrokontrolerów różnych serii. \\ \hline
\textbf{Python: \texttt{cantools}} & \textbf{40.7.1} & \textbf{Generator kodu C z plików DBC.} \\ \hline
Python: \texttt{python-can} & (auto) & Instalowana automatycznie jako zależność \texttt{cantools}. \\ \hline
\end{tabular}
\caption{Wymagania systemowe projektu}
\label{tab:wymagania}
\end{table}

\subsection{Instalacja zależności Python}

Skrypt generujący kod (\texttt{dbc/setup.py}) jest uruchamiany automatycznie przez CMake i wymaga zainstalowanych bibliotek w środowisku Pythona. Biblioteka \texttt{cantools} w wersji 40.7.1 pociąga za sobą wymagane zależności (m.in. \texttt{python-can}, \texttt{bitstruct}).

Aby zainstalować wymagane pakiety, zaleca się użycie poniższego polecenia w terminalu:

\begin{lstlisting}[style=bashstyle, caption=Instalacja wymaganych bibliotek Python]
pip install cantools==40.7.1
\end{lstlisting}

Brak tych bibliotek spowoduje błąd konfiguracji CMake na etapie \texttt{Checking and regenerating DBC files}.

\newpage

\section{Integracja i Konfiguracja w CMake}

Biblioteka jest dystrybuowana jako moduł CMake. Aby użyć jej w projekcie, należy zmodyfikować główny plik \texttt{CMakeLists.txt} projektu.

\subsection{Dodawanie biblioteki do projektu}

W głównym pliku \texttt{CMakeLists.txt} (np. w katalogu głównym projektu STM32 lub węzła ROS2) należy dodać podkatalog z biblioteką oraz zlinkować ją z plikiem wykonywalnym.

\begin{lstlisting}[style=cmakestyle, caption=Podstawowa konfiguracja CMake]
# 1. Dodanie podkatalogu biblioteki
add_subdirectory(Core/Inc/PUTM_EV_CAN_LIBRARY)

# 2. Linkowanie biblioteki z twoim projektem
# 'twoj_projekt' to nazwa zdefiniowana w add_executable()
target_link_libraries(twoj_projekt PRIVATE putm_ev_can)
\end{lstlisting}

\subsection{Wybór konfiguracji (STM32 vs ROS2)}

Biblioteka wspiera dwa backhendy sprzętowe. Wybór następuje poprzez ustawienie zmiennej \texttt{PUTM\_BACKEND} \textbf{przed} wywołaniem \texttt{add\_subdirectory}.

\subsubsection{Konfiguracja dla STM32}
Jest to domyślna konfiguracja. Wymaga ona, aby w projekcie dostępna była biblioteka HAL.

\begin{lstlisting}[style=cmakestyle, caption=Konfiguracja dla STM32]
set(PUTM_BACKEND "STM32" CACHE STRING "Select backend" FORCE)
add_subdirectory(Core/Inc/PUTM_EV_CAN_LIBRARY)
\end{lstlisting}

\subsubsection{Konfiguracja dla ROS2 / Linux}
W tym trybie biblioteka korzysta z systemowego SocketCAN.

\begin{lstlisting}[style=cmakestyle, caption=Konfiguracja dla Linux/ROS2]
set(PUTM_BACKEND "LINUX" CACHE STRING "Select backend" FORCE)
add_subdirectory(lib/PUTM_EV_CAN_LIBRARY)
\end{lstlisting}

\subsection{Włączenie obsługi CAN FD}
Domyślnie biblioteka pracuje w trybie Classic CAN. Aby włączyć obsługę ramek do 64 bajtów:

\begin{lstlisting}[style=cmakestyle]
option(PUTM_USE_CAN_FD "Enable CAN FD" ON)
\end{lstlisting}
\newpage
\section{Konwencja Nomenklatury Magistral}

W projekcie pojazdu zdefiniowano trzy fizyczne magistrale CAN. Aby zachować spójność w kodzie (na STM32 i w ROS2), \textbf{należy bezwzględnie} stosować poniższe nazewnictwo instancji klasy \texttt{AppCAN}.

\begin{table}[H]
\centering
\begin{tabular}{@{}lll@{}}
\toprule
\textbf{Nazwa Instancji} & \textbf{Funkcja Magistrali} & \textbf{Opis} \\ \midrule
\texttt{\textbf{can\_m}} & \textbf{Main} & Główna magistrala, większość danych, telemetria. \\
\texttt{\textbf{can\_pt}} & \textbf{Powertrain (ECU/INV)} & Komunikacja wyłącznie falownik - silniki. \\
\texttt{\textbf{can\_dv}} & \textbf{Driverless} & Wymiana sygnałów w trybie jazdy autonomicznej. \\ \bottomrule
\end{tabular}
\caption{Obowiązująca nomenklatura instancji AppCAN}
\label{tab:nomenklatura}
\end{table}

Wszystkie przykłady w dalszej części dokumentacji będą odwoływać się do tych konkretnych nazw zmiennych.

\section{Poprawna inicjalizacja}

Biblioteka nie używa dynamicznej alokacji pamięci. Instancje muszą być zadeklarowane jako zmienne globalne.

\subsection{Krok 1: Instancje Globalne}
W pliku \texttt{main.cpp} (lub odpowiedniku w ROS2) należy utworzyć obiekty dla wykorzystywanych magistral.

\begin{lstlisting}
/* Private variables */
// Definicja instancji zgodnie z przyjeta nomenklatura
AppCAN can_m;   // Main Bus
AppCAN can_pt;  // Powertrain Bus
AppCAN can_dv;  // Driverless Bus
\end{lstlisting}

\subsection{Krok 2: Wywołanie Init}
Inicjalizację należy przeprowadzić wewnątrz funkcji \texttt{main()}, \textbf{po} konfiguracji peryferiów przez CubeMX. Każdej instancji przypisujemy odpowiedni uchwyt sprzętowy (\texttt{hfdcanX} lub \texttt{hcanX}).

\begin{lstlisting}
/* USER CODE BEGIN 2 */

// 1. Inicjalizacja Main CAN (np. FDCAN1)
if (!can_m.Init(&hfdcan1)) {
    Error_Handler(); // Blad inicjalizacji Main
}

// 2. Inicjalizacja Powertrain CAN (np. FDCAN2)
if (!can_pt.Init(&hfdcan2)) {
    Error_Handler(); // Blad inicjalizacji PT
}
/* USER CODE END 2 */
\end{lstlisting}

\textbf{Uwaga:} Biblioteka automatycznie wykrywa rodzinę procesora (np. G4 vs F4) i dostosowuje typ uchwytu (FDCAN vs bxCAN).

\section{Diagnostyka LED}

Biblioteka posiada wbudowany mechanizm automatycznej sygnalizacji stanu magistrali za pomocą diody LED. Funkcjonalność ta pozwala na szybką diagnostykę problemów (np. odpięta wtyczka, uszkodzona magistrala, błędy ramek) bez konieczności podłączania debuggera.

\subsection{Konfiguracja Diody}
Dla każdej instancji \texttt{AppCAN} (np. \texttt{can\_m}, \texttt{can\_pt}) można przypisać osobną diodę. Konfigurację należy wykonać \textbf{raz}, zaraz po pomyślnej inicjalizacji biblioteki.

\begin{lstlisting}
if (can_m.Init(&hfdcan1)) {
    // Przypisanie diody (np. LED1) do tej magistrali
    can_m.ConfigStatusLed(LED1_GPIO_Port, LED1_Pin);
}
\end{lstlisting}

\subsection{Znaczenie Sygnałów (Mruganie)}
Po skonfigurowaniu, biblioteka automatycznie steruje diodą wewnątrz metody \texttt{Poll()}. Użytkownik nie musi (i nie powinien) ręcznie sterować tymi pinami.

\begin{table}[H]
\centering
\begin{tabular}{@{}lll@{}}
\toprule
\textbf{Zachowanie Diody} & \textbf{Stan Magistrali} & \textbf{Opis} \\ \midrule
\textbf{Mruganie 1 Hz} & \textbf{OK} & Magistrala działa poprawnie (Error Active). \\
(500ms ON, 500ms OFF) & & Komunikacja jest możliwa. \\ \midrule
\textbf{Nierówne Mruganie} & \textbf{WARNING} & Wykryto błędy transmisji (Error Passive / Warning). \\
(Podwójny błysk co 1s) & & Może oznaczać zakłócenia lub błędy ramek. \\ \midrule
\textbf{Świeci Ciągle} & \textbf{BUS OFF} & Krytyczny błąd magistrali. \\
(Solid ON) & & Kontroler CAN odłączył się od sieci (np. zwarcie). \\ \bottomrule
\end{tabular}
\caption{Kody diagnostyczne diody statusowej}
\label{tab:led_codes}
\end{table}


\newpage
\section{Wysyłanie danych (TX)}

Wysyłanie odbywa się poprzez metodę \texttt{Send} wywołaną na odpowiedniej instancji magistrali (zazwyczaj \texttt{can\_m}). Pamiętaj o używaniu poprawnego prefiksu \texttt{PUTM\_CAN\_M\_} (lub innego, zależnie od pliku DBC).

\begin{lstlisting}
// Przyklad w petli glownej (co 100ms)
if (HAL_GetTick() - last_tx > 100) {
    last_tx = HAL_GetTick();

    // 1. Tworzenie i zerowanie struktury
    PUTM_CAN_M_pc_main_data_t tx_msg = {0};

    // 2. Wypelnianie danych
    tx_msg.vehicle_speed = 125;
    tx_msg.inverters_ready = 1;
    
    // 3. Wyslanie na magistrale MAIN (can_m)
    if (can_m.Send(PUTM_CAN_M_PC_MAIN_DATA_FRAME_ID, tx_msg)) {
        // Opcjonalnie: dodatkowa logika po wyslaniu
    }
}
\end{lstlisting}

\newpage
\section{Odbieranie danych (RX)}

Odbiór w nowej wersji biblioteki działa w trybie \textbf{Interrupt-Driven (Przerwania)}. Oznacza to, że dane są pobierane natychmiast po ich nadejściu, a nie w pętli głównej.

\subsection{Krok 1: Rejestracja Callbacku}
Rejestrację wykonujemy raz, przed pętlą główną.
\textbf{OSTRZEŻENIE:} Funkcja lambda (callback) jest wywoływana w kontekście przerwania sprzętowego (ISR). Należy w niej unikać:
\begin{itemize}
    \item Funkcji blokujących (\texttt{HAL\_Delay}).
    \item Długich operacji matematycznych.
    \item Funkcji wejścia/wyjścia (np. \texttt{printf}).
\end{itemize}
Zaleca się jedynie ustawianie flag, zmiennych \texttt{volatile} lub szybkie sterowanie pinami.

\begin{lstlisting}
// Odbior danych z Main CAN
can_m.RegisterCallback<PUTM_CAN_M_rearbox_safety_t>(
    PUTM_CAN_M_REARBOX_SAFETY_FRAME_ID,
    [](const PUTM_CAN_M_rearbox_safety_t& msg) {
        
        // Ta funkcja wykonuje sie w PRZERWANIU!
        // Przyklad bezpiecznej operacji: sterowanie LED
        if (msg.safety_tsmp == 1) {
             HAL_GPIO_WritePin(LED3_GPIO_Port, LED3_Pin, GPIO_PIN_RESET);
        } else {
             HAL_GPIO_WritePin(LED3_GPIO_Port, LED3_Pin, GPIO_PIN_SET);
        }
    }
);
\end{lstlisting}

\subsection{Krok 2: Polling w pętli głównej}
Mimo że odbiór danych jest asynchroniczny, metoda \texttt{Poll()} nadal musi być wywoływana. Jej zadaniem jest teraz obsługa \textbf{diody diagnostycznej}.

\begin{lstlisting}
while (1) {
    // Obsluga mrugania dioda diagnostyczna
    // (Odbior ramek dzieje sie automatycznie w tle)
    can_m.Poll();
    can_pt.Poll();
    can_dv.Poll();
    
    // ... reszta kodu ...
}
\end{lstlisting}

\section{Najczęściej popełniane błędy}

\begin{itemize}
    \item \textbf{Użycie złego prefiksu:} Biblioteka teraz używa \texttt{PUTM\_CAN\_M\_} (dla Main DBC). Użycie starego \texttt{PUTM\_CAN\_1\_} spowoduje błąd kompilacji.
    \item \textbf{Brak słowa \texttt{volatile}:} Jeśli zmienna jest modyfikowana w callbacku (RX) i czytana w \texttt{main()}, musi być oznaczona jako \texttt{volatile}, inaczej kompilator może ją źle zoptymalizować.
    \item \textbf{Pomylenie instancji:} Wysłanie wiadomości na \texttt{can\_m} zamiast \texttt{can\_pt}, przez co falownik jej nie otrzyma.
    \item \textbf{Brak wywołania \texttt{Poll()}:} Skutkuje brakiem sygnalizacji LED (dioda nie będzie mrugać mimo poprawnej pracy).
\end{itemize}


\newpage
\section{Szczegółowa Architektura Systemu}

Biblioteka została zaprojektowana w oparciu o wzorzec projektowy \textit{Fasada} (Facade) oraz \textit{Wstrzykiwanie Zależności}. Architektura jest podzielona na cztery główne warstwy.

\begin{figure}[H]
    \centering
    % Miejsce na diagram (np. wygenerowany z Mermaid)
    % \includegraphics[width=0.9\textwidth]{architecture.png}
    \caption{Diagram warstwowy biblioteki}
\end{figure}

\subsection{Warstwa 1: Fasada Aplikacji (\texttt{app\_can\_lib})}
Jest to jedyna warstwa, z którą bezpośrednio intereaguje użytkownik końcowy.
\begin{itemize}
    \item \textbf{Odpowiedzialność:} Ukrywa złożoność wewnętrzną, zarządza czasem życia obiektów i obsługuje diagnostykę LED w metodzie \texttt{Poll()}.
\end{itemize}

\subsection{Warstwa 2: Logika i Routing}
\subsubsection{Interfejs (\texttt{can\_interface})}
Rejestruje callback odbiorczy w warstwie HAL. Gdy przyjdzie przerwanie, HAL przekazuje surowe dane tutaj, a interfejs pakuje je w struktury C++ (korzystając z DBC).

\subsubsection{Router (\texttt{message\_handler})}
Przechowuje listę callbacków użytkownika. Zapewnia bezpieczeństwo wątkowe przy rejestracji nowych funkcji.

\subsection{Warstwa 3: Abstrakcja Sprzętowa HAL (\texttt{can\_hal})}
Warstwa izolująca kod od sprzętu. Została znacznie rozbudowana w nowej wersji.

\begin{itemize}
    \item \textbf{Selektor (\texttt{stm32\_hal\_selector.hpp}):} Automatycznie wykrywa rodzinę procesora (np. G484xx) i dołącza odpowiednie nagłówki HAL. Definiuje typ \texttt{CanHandleType} (FDCAN dla G4/H7, bxCAN dla F4/L4).
    \item \textbf{Dispatcher Przerwań:} Klasa \texttt{Stm32CanHal} posiada statyczną tablicę, która mapuje uchwyty sprzętowe na obiekty C++. Pozwala to globalnej funkcji przerwania C (\texttt{HAL\_FDCAN\_RxFifo0Callback}) odnaleźć właściwą instancję biblioteki i przekazać jej dane.
\end{itemize}

\subsection{Warstwa 4: Definicje Danych (\texttt{dbc/generated})}
Kod generowany automatycznie przez skrypt Python na podstawie plików \texttt{.dbc}. Zawiera definicje struktur (np. \texttt{PUTM\_CAN\_M\_bms\_lv\_main\_t}) oraz funkcje pakujące.

\end{document}